\section*{Problem 388. Equal products of two disjoint blocks of consecutive integers}

\subsection*{1) ROUND-2 OBJECTIVE}

Path (C): obstruction/correction.  Round~1 identified a real ambiguity in the problem statement: one can ask either

\begin{itemize}
\item for each fixed pair of lengths $(k_1,k_2)$, are there only finitely many solutions?, or
\item over all varying lengths and shifts, are there only finitely many solutions?
\end{itemize}

The most promising Round-2 direction is therefore to separate these two meanings rigorously.  The key new fact is that the fixed-length problem is already a theorem, whereas the all-parameters problem remains open.

\subsection*{2) Round-1 FOUNDATION USED}

I use the following Round-1 results without reproving them.

\begin{itemize}
\item \textbf{Round-1 Lemma 1.}  Under the disjointness condition $m_1+k_1\le m_2$, any solution of
\[
\prod_{i=1}^{k_1}(m_1+i)=\prod_{j=1}^{k_2}(m_2+j)
\]
must satisfy $k_1>k_2$.

\item \textbf{Round-1 Lemma 2.}  If $A=m_1+k_1$, then every prime factor of every integer in the later block $\{m_2+1,\dots,m_2+k_2\}$ is at most $A$.
\end{itemize}

Only Lemma~1 is logically needed for the main correction below; Lemma~2 remains a useful structural obstruction for the still-open global version.

\subsection*{3) NEW INSIGHT / TOOL (ROUND-2)}

The genuinely new ingredient is that the fixed-length interpretation is not open.

\begin{itemize}
\item For each fixed pair of lengths, finiteness follows from a theorem of Beukers, Shorey and Tijdeman on equations of the shape
\[
x(x+1)\cdots(x+k-1)=y(y+1)\cdots(y+k+l-1)
\]
with fixed $k,l$.

\item More generally, for each fixed rational ratio of the two lengths, finiteness follows from a theorem of Saradha, Shorey and Tijdeman on
\[
d_1\,x(x+1)\cdots(x+lk-1)=d_2\,y(y+1)\cdots(y+mk-1).
\]

\item In the low-ratio range $k_1/k_2\in\{2,3,4,5,6\}$ there are in fact no solutions satisfying the present constraint $k_1,k_2>3$.
\end{itemize}

So the Round-1 ambiguity can now be resolved sharply:
the fixed-length reading is settled; the all-parameters reading is still open.

\subsection*{4) ATTACK PLAN (ROUND-2)}

The exact Round-1 gap was this: no theorem had been supplied converting the present equation into a known finiteness result for fixed lengths.

I now do three things.

\begin{enumerate}
\item Convert the present notation to the standard ``shorter block versus longer block'' form used in the literature, and apply the fixed-$(k,l)$ theorem.

\item Strengthen this to fixed \emph{length ratio} using the Saradha--Shorey--Tijdeman theorem.

\item Record one further sharp consequence: if the longer length is $m$ times the shorter one with $2\le m\le 6$, then there are no solutions with both lengths $>3$.
\end{enumerate}

This closes the fixed-length gap completely and isolates the true open remainder.

\subsection*{5) WORK (ROUND-2)}

We study
\[
\prod_{i=1}^{k_1}(m_1+i)=\prod_{j=1}^{k_2}(m_2+j),
\qquad
k_1,k_2>3,\qquad m_1+k_1\le m_2.
\]

By Round-1 Lemma~1, any such solution satisfies $k_1>k_2$.

\medskip
\noindent\textbf{Claim 1.}
\emph{For each fixed pair $(k_1,k_2)$ with $k_1>k_2\ge 4$, there are only finitely many integer solutions $(m_1,m_2)$ satisfying the equation and the disjointness condition.}

\medskip
\noindent\textbf{Proof.}
Fix $k_1>k_2\ge 4$ and put
\[
k:=k_2,\qquad l:=k_1-k_2,\qquad x:=m_2+1,\qquad y:=m_1+1.
\]
Then the equation becomes
\[
x(x+1)\cdots(x+k-1)=y(y+1)\cdots(y+k+l-1).
\]
Since $m_1+k_1\le m_2$, we have
\[
x=m_2+1\ge m_1+k_1+1=y+k+l.
\]
In particular $x$ and $y$ are positive integers, and we are exactly in the fixed-$(k,l)$ situation treated by Beukers--Shorey--Tijdeman.

Their theorem says that for fixed positive integers $k,l$, the equation
\[
x(x+1)\cdots(x+k-1)=y(y+1)\cdots(y+k+l-1)
\]
has only finitely many positive integer solutions $(x,y)$.

Therefore, for each fixed pair $(k_1,k_2)$, the original problem has only finitely many solutions $(m_1,m_2)$.
\hfill\textit{QED.}

\medskip
\noindent\textbf{Thus the fixed-length interpretation is fully settled.}

\medskip
\noindent\textbf{Claim 2.}
\emph{Fix coprime positive integers $l<m$.  Then there are only finitely many solutions whose two lengths are in the ratio $l:m$, i.e. whose lengths are $(lk,mk)$ for some positive integer $k$.}

\medskip
\noindent\textbf{Proof.}
Write the shorter length as $lk$ and the longer one as $mk$, with $(l,m)=1$ and $l<m$.
After the same change of variables as above, the equation becomes
\[
x(x+1)\cdots(x+lk-1)=y(y+1)\cdots(y+mk-1).
\]
This is the special case $d_1=d_2=1$ of the Saradha--Shorey--Tijdeman theorem on
\[
d_1\,x(x+1)\cdots(x+lk-1)=d_2\,y(y+1)\cdots(y+mk-1).
\]
That theorem states that either $\max(x,y,k)$ is bounded by a constant depending only on $m,d_1,d_2$, or one falls into an exceptional case
\[
m=2,\qquad k=2,\qquad d_1=2d_2^2,
\]
together with an auxiliary square condition.  In our specialization $d_1=d_2=1$, the exceptional condition $d_1=2d_2^2$ becomes $1=2$, impossible.

Hence $\max(x,y,k)$ is bounded, and therefore only finitely many solutions exist for each fixed ratio $l:m$.
\hfill\textit{QED.}

\medskip
\noindent\textbf{Claim 3.}
\emph{If the longer block has length $m$ times the shorter one with $2\le m\le 6$, then there is no solution satisfying the present hypothesis $k_1,k_2>3$.}

\medskip
\noindent\textbf{Proof.}
Saradha--Shorey and Mignotte--Shorey proved that for $2\le m\le 6$ the only solution to
\[
x(x+1)\cdots(x+k-1)=y(y+1)\cdots(y+mk-1)
\]
in positive integers is
\[
x=8,\qquad y=1,\qquad k=3,\qquad m=2.
\]
Here the two block lengths are $3$ and $6$.

Under our problem's standing condition $k_1,k_2>3$, the shorter length is at least $4$, so this exceptional solution is excluded.  Therefore no solutions exist in the ratio range $k_1/k_2\in\{2,3,4,5,6\}$.
\hfill\textit{QED.}

\medskip
\noindent\textbf{Claim 4.}
\emph{For every fixed $N$, there are only finitely many solutions with $\max(k_1,k_2)\le N$ and $k_2\nmid 2k_1$.}

\medskip
\noindent\textbf{Proof.}
By Round-1 Lemma~1, the shorter length is $k_2$ and the longer length is $k_1$.
A theorem of Hajdu and Tijdeman implies that for every fixed $N$, there are only finitely many pairs of disjoint blocks of positive integers of sizes at most $N$ with equal products whenever the smaller size does not divide twice the larger size.

Our problem is a special case of that theorem, because our two blocks are themselves disjoint consecutive blocks.
Thus, for fixed $N$, there are only finitely many solutions with $\max(k_1,k_2)\le N$ and $k_2\nmid 2k_1$.
\hfill\textit{QED.}

\subsection*{6) ADVERSARIAL VERIFICATION}

I check the possible failure points.

\begin{itemize}
\item \textbf{Variable conversion.}
The fixed-length theorem is stated for a shorter block of length $k$ and a longer block of length $k+l$.
Round-1 Lemma~1 guarantees exactly that ordering here, so the conversion is legitimate.

\item \textbf{Disjointness.}
The literature theorems do not require the extra condition $x\ge y+k+l$; they only need positive integer variables.
Hence adding disjointness can only restrict the solution set, not invalidate the theorem.

\item \textbf{Exceptional case in the ratio theorem.}
The only exceptional branch of the Saradha--Shorey--Tijdeman theorem requires $d_1=2d_2^2$.
In our application $d_1=d_2=1$, so the exceptional branch is impossible.

\item \textbf{Low-ratio theorem versus the present hypotheses.}
The unique solution for ratios $2\le m\le 6$ has smaller length $3$, which violates the standing assumption $k_1,k_2>3$.
So the conclusion ``no solutions'' is correct for the present problem.

\item \textbf{Did Round-1 overstate openness?}
Yes, for the fixed-length reading.
Round~1 treated that case as unresolved, but the above theorems show that interpretation is already settled.
The all-parameters reading is the truly open one.
\end{itemize}

\subsection*{7) FINAL}

\textbf{UNRESOLVED (BUT STRICTLY ADVANCED).}

The strict advance is the following precise correction.

\begin{enumerate}
\item The \emph{fixed-length} interpretation of Problem~388 is already true and can be proved rigorously from Beukers--Shorey--Tijdeman: for each fixed pair $(k_1,k_2)$ with $k_1>k_2\ge 4$, only finitely many disjoint-block solutions exist.

\item More generally, for each fixed rational ratio of the two lengths, only finitely many solutions exist.

\item If the ratio of lengths is an integer $m$ with $2\le m\le 6$, then under the standing hypothesis $k_1,k_2>3$ there are no solutions at all.

\item What remains open is the \emph{all-parameters} problem: are there only finitely many quadruples $(m_1,k_1,m_2,k_2)$ with
\[
k_1,k_2>3,\qquad m_1+k_1\le m_2,\qquad
\prod_{i=1}^{k_1}(m_1+i)=\prod_{j=1}^{k_2}(m_2+j)?
\]
Round-1 Lemma~2 still gives the main structural obstruction here: the later block must be exceptionally smooth.
\end{enumerate}

\subsection*{8) COMPLETION ESTIMATE (MANDATORY)}

COMPLETION: 80\%

\subsection*{9) REFERENCES}

\begin{enumerate}
\item F. Beukers, T. N. Shorey, and R. Tijdeman, results on the equation
\[
x(x+1)\cdots(x+k-1)=y(y+1)\cdots(y+k+l-1)
\]
for fixed $k,l$; summarized in the survey ``Highlights in the Research Work of T. N. Shorey.''

\item S. Saradha and T. N. Shorey, and A. Mignotte and T. N. Shorey, results on
\[
x(x+1)\cdots(x+k-1)=y(y+1)\cdots(y+mk-1)
\]
for $2\le m\le 6$; summarized in the same survey.

\item S. Saradha, T. N. Shorey, and R. Tijdeman, results on
\[
d_1\,x(x+1)\cdots(x+lk-1)=d_2\,y(y+1)\cdots(y+mk-1).
\]

\item L. Hajdu and R. Tijdeman, \emph{The Diophantine equation $f(x)=g(y)$ for polynomials with simple rational roots}, especially Theorem~10.1.
\end{enumerate}



